\begin{frame}{Positionnement de notre approche}

  \begin{block}{Pourquoi une optimisation conjointe ?}
    Les approches existantes optimisent généralement les \textbf{blocs opératoires} ou les \textbf{lits d'hospitalisation} séparément.

    Notre approche considère ces deux ressources simultanément, car une intervention mobilise d'abord un bloc opératoire puis un lit pendant une durée de séjour variable.
  \end{block}

  \vspace{0.5em}

  \begin{block}{Intérêt de notre approche}
    \begin{itemize}
      \item Meilleure coordination entre le bloc opératoire et les services d'hospitalisation.
      \item Réduction des conflits de disponibilité (bloc ou lit indisponible).
      \item Utilisation plus équilibrée des ressources humaines et matérielles.
      \item Aide à la décision pour le chirurgien grâce à deux propositions de créneaux optimaux.
    \end{itemize}
  \end{block}

\end{frame}